% =============================================================================
% 从地球到星辰
% 星际旅行、恒星际航行与外星文明接触的数学框架
%
% 作者：GitHub Copilot
%       人工智能研究部门
%       GitHub | 微软公司
% 日期：2026年4月4日
% =============================================================================
%
% 编译命令：xelatex paper_zh.tex
%           bibtex paper_zh
%           xelatex paper_zh.tex
%           xelatex paper_zh.tex
% =============================================================================

\documentclass[12pt,a4paper]{article}

% --- 中文支持（需要 XeLaTeX）-------------------------------------------------
\usepackage{ctex}

% --- 数学与符号 ---------------------------------------------------------------
\usepackage{amsmath}
\usepackage{amssymb}
\usepackage{mathtools}
\usepackage{bm}

% --- 页面布局 -----------------------------------------------------------------
\usepackage[a4paper,
            top=3cm, bottom=3cm,
            left=2.8cm, right=2.8cm,
            headheight=15pt]{geometry}

% --- 图形与颜色 ---------------------------------------------------------------
\usepackage{graphicx}
\usepackage{xcolor}
\usepackage{tikz}
\usepackage{pgfplots}
\pgfplotsset{compat=1.18}
\usetikzlibrary{arrows.meta, calc, decorations.markings}

% --- 表格 ---------------------------------------------------------------------
\usepackage{booktabs}
\usepackage{array}
\usepackage{tabularx}
\usepackage{multirow}

% --- 浮动体与标题 -------------------------------------------------------------
\usepackage{float}
\usepackage[labelfont=bf, font=small]{caption}
\usepackage{subcaption}

% --- 列表 ---------------------------------------------------------------------
\usepackage{enumitem}

% --- 单位 ---------------------------------------------------------------------
\usepackage{siunitx}

% --- 参考文献与超链接 ---------------------------------------------------------
\usepackage[authoryear, round, sort]{natbib}
\usepackage[colorlinks=true,
            linkcolor=blue!65!black,
            citecolor=green!50!black,
            urlcolor=blue!70!black,
            pdftitle={从地球到星辰},
            pdfauthor={GitHub Copilot}
           ]{hyperref}

% --- 其他排版 -----------------------------------------------------------------
\usepackage{setspace}
\usepackage{fancyhdr}

% =============================================================================
% 自定义命令
% =============================================================================
\newcommand{\dd}{\mathrm{d}}
\newcommand{\ee}{\mathrm{e}}
\newcommand{\kB}{k_{\mathrm{B}}}
\newcommand{\AU}{\,\text{AU}}
\newcommand{\ly}{\,\text{光年}}
\newcommand{\Msun}{M_{\odot}}
\DeclareMathOperator{\atanh}{atanh}

% =============================================================================
% 页面样式
% =============================================================================
\pagestyle{fancy}
\fancyhf{}
\fancyhead[L]{\small\textsl{GitHub Copilot~——~\textit{从地球到星辰}}}
\fancyhead[R]{\small\textsl{2026年4月4日}}
\fancyfoot[C]{\small\thepage}
\renewcommand{\headrulewidth}{0.4pt}
\setlength{\parskip}{5pt}
\setlength{\parindent}{2em}

% =============================================================================
% 标题
% =============================================================================
\title{%
  \vspace{-0.8cm}%
  {\rule{\linewidth}{1.5pt}}\\[0.5em]
  {\LARGE\bfseries 从地球到星辰}\\[0.4em]
  {\large\itshape 星际旅行、恒星际航行与\\
   外星文明接触的数学框架}\\[0.5em]
  {\rule{\linewidth}{0.8pt}}%
}
\author{%
  \textbf{GitHub Copilot}\\[0.3em]
  \small 人工智能研究部门\\
  \small GitHub~|~微软公司\\[0.3em]
  \small\texttt{copilot@github.com}%
}
\date{\small 2026年4月4日}

% =============================================================================
\begin{document}
% =============================================================================

\maketitle
\thispagestyle{empty}

% --- 摘要 --------------------------------------------------------------------
\begin{abstract}
\noindent
本文对人类迈向行星际定居和恒星际探索的长期轨迹进行了严格的数学论证。从轨道转移的经典力学与齐奥尔科夫斯基火箭方程出发，逐步推进至先进推进概念——电推进、核推进和反物质推进——直至相对论性太空飞行领域，在那里洛伦兹因子成为任务规划的核心。随后，通过对德雷克方程的贝叶斯重新表述，对费米悖论进行概率框架下的量化，并分析星际通信的工程与信息论要求。一个原创章节提出了首次接触的形式化数学语言，借助数论、普适物理常数和香农信息论构建体系。论文最后探讨星际关系建立的博弈论模型，论证数学推理——普适而独立于文化——是宇宙外交的唯一可行基础。所有结论均辅以推导过程和适用于当前及近未来技术的数值估算。

\bigskip\noindent
\textbf{关键词：}轨道力学，星际推进，相对论性太空飞行，德雷克方程，SETI，星际通信，首次接触，博弈论，费米悖论。
\end{abstract}

\tableofcontents
\clearpage

% =============================================================================
\section{引言}
% =============================================================================

人类自古以来便仰望星空，怀揣着敬畏与探索的渴望。从神话传说与航海定向——北极星引领水手归家，昴星团标示农时——到哥白尼、开普勒、伽利略与牛顿，天文学逐步演变为一门严谨的数学科学。1957年10月4日，人造卫星"斯普特尼克"成为首个进入地球轨道的人造物体，将梦想化为工程现实。此后，机器人探测器相继造访了太阳系所有行星，火星车以极高精度绘制火星地形图，而双"旅行者"号探测器历经数十年飞行，已穿越日球层顶，进入星际空间。

然而，将我们与哪怕最近的恒星邻居分隔开来的距离，以任何人类尺度衡量都是令人咋舌的。半人马座$\alpha$星系位于 \SI{4.37}{\ly} 之外——约 \SI{4.13e13}{\kilo\meter}。以旅行者1号约 \SI{17}{\kilo\meter\per\second} 的速度，这段旅程需要大约 \num{73000} 年。行星际旅行困难重重；恒星际旅行则几乎超出人类的想象。然而物理学并不禁止它。力学、热力学和电磁学定律——以巨大但有限的代价——允许向其他恒星发送载人或无人探测器，并使其在人类有生之年内抵达目的地。

本文旨在构建这一可能性的连贯数学图景，以及它所蕴含的文明意义。第2节介绍轨道运动的经典力学；第3节引入火箭方程；第4节将上述基础转化为太阳系内的具体任务设计；第5节概述先进推进技术；第6节将狭义相对论引入舞台中心；第7--9节讨论我们是否孤独、如何跨越恒星距离通信，以及哪些数学结构可作为通用语言；最后几节审视星际接触的逻辑、伦理与外交框架。

% =============================================================================
\section{天体力学：运动的数学语言}
% =============================================================================

\subsection{牛顿万有引力定律}

天体力学的基础是牛顿平方反比定律：
\begin{equation}
  \bm{F} = -\frac{GM_1 M_2}{r^2}\hat{\bm{r}},
\end{equation}
其中 $G = \SI{6.674e-11}{\meter\cubed\per\kilo\gram\per\second\squared}$。
对于质量 $m \ll M_\odot$ 的测试质量绕太阳运动，极坐标 $(r,\theta)$ 下的运动方程为：
\begin{equation}
  \ddot{r} - r\dot{\theta}^2 = -\frac{GM_\odot}{r^2}, \qquad
  \frac{\dd}{\dd t}(r^2\dot{\theta}) = 0.
\end{equation}
第二个方程体现了比角动量 $h = r^2\dot{\theta} = \text{常数}$ 的守恒，即开普勒第二定律的微分形式。

\subsection{开普勒定律}

轨道的极坐标方程为：
\begin{equation}
  r(\theta) = \frac{a(1-e^2)}{1 + e\cos\theta},
\end{equation}
其中 $a$ 为半长轴，$e$ 为轨道离心率。开普勒三大定律由此直接推出：
\begin{enumerate}[label=\textbf{K\arabic*.}]
  \item 行星轨道为椭圆，太阳位于其中一个焦点。
  \item 径矢在相等时间内扫过相等面积。
  \item 轨道周期的平方与半长轴的立方成正比：
    \begin{equation}
      T^2 = \frac{4\pi^2}{GM_\odot}\,a^3.
    \end{equation}
\end{enumerate}

\subsection{活力公式}

单位质量的总机械能守恒：
\begin{equation}
  \varepsilon = \frac{v^2}{2} - \frac{GM_\odot}{r} = -\frac{GM_\odot}{2a},
\end{equation}
整理后得到活力公式（vis-viva）：
\begin{equation}
  \boxed{v^2 = GM_\odot\!\left(\frac{2}{r} - \frac{1}{a}\right).}
\end{equation}
在地球平均轨道半径处，$v_\oplus \approx \SI{29.78}{\kilo\meter\per\second}$。

\subsection{逃逸速度}

令总能量为零，得逃逸速度：
\begin{equation}
  v_{\mathrm{逃逸}}(r) = \sqrt{\frac{2GM_\odot}{r}} = \sqrt{2}\,v_{\mathrm{圆轨道}}(r).
\end{equation}
从地球表面逃逸的速度为 $v_{\mathrm{逃逸},\oplus} \approx \SI{11.19}{\kilo\meter\per\second}$。

% =============================================================================
\section{火箭方程：动能的代价}
% =============================================================================

\subsection{从牛顿第二定律的推导}

对变质量系统应用牛顿第二定律：
\begin{equation}
  m\frac{\dd v}{\dd t} = -v_e\frac{\dd m}{\dd t},
\end{equation}
其中 $v_e$ 为推进剂在火箭参考系中的喷射速度。从初始质量 $m_0$ 积分到末质量 $m_f$：
\begin{equation}
  \boxed{\Delta v = v_e \ln\!\frac{m_0}{m_f} = I_{\mathrm{sp}}\,g_0\,\ln\!\frac{m_0}{m_f}.}
  \label{eq:tsiolkovsky_zh}
\end{equation}
这个方程由齐奥尔科夫斯基于1903年首先发表，是火箭工程的基石~\citep{Tsiolkovsky1903}。比冲 $I_{\mathrm{sp}} = v_e/g_0$（单位：秒）是任何推进系统性能的核心指标。

\subsection{质量比问题}

给定 $\Delta v$ 所需的质量比可由式~\eqref{eq:tsiolkovsky_zh} 直接得到：
\begin{equation}
  \frac{m_0}{m_f} = \exp\!\left(\frac{\Delta v}{v_e}\right).
\end{equation}
这种指数依赖对高 $\Delta v$ 任务而言是灾难性的。表~\ref{tab:dv_zh}
以氢氧发动机（$I_{\mathrm{sp}} = \SI{450}{\second}$，
$v_e = \SI{4.41}{\kilo\meter\per\second}$）为例说明了这一问题。

\begin{table}[h]
  \centering
  \caption{代表性任务的 $\Delta v$ 预算与质量比
           （LOX/LH$_2$ 推进，$I_{\mathrm{sp}} = 450$\,s）。}
  \label{tab:dv_zh}
  \begin{tabularx}{\textwidth}{Xrr}
    \toprule
    任务 & $\Delta v$ (km\,s$^{-1}$) & 质量比 $m_0/m_f$ \\
    \midrule
    低地球轨道 & 9.5  & 8.6   \\
    地月转移轨道注入 & 12.5 & 17.0 \\
    地火转移轨道注入 & 11.6 & 13.9 \\
    直飞木星 & 20.0 & 93    \\
    从地球逃逸太阳系 & 16.6 & 43    \\
    $\Delta v = 1\,000$\,km\,s$^{-1}$（$0.33\%\,c$） &
    \num{1000} & $\sim 10^{98}$ \\
    \bottomrule
  \end{tabularx}
\end{table}

最后一行明确无误地表明：仅凭火箭方程，化学推进无法达到恒星际速度；
必须采用替代性的推进方案。

% =============================================================================
\section{行星际任务设计}
% =============================================================================

\subsection{霍曼转移轨道}

在半径分别为 $r_1$ 和 $r_2$ 的共面圆轨道之间能量最优的两次脉冲转移，
即霍曼转移（1925年）~\citep{Hohmann1925}。所需速度增量为：
\begin{align}
  \Delta v_1 &= \sqrt{\frac{GM_\odot}{r_1}}
               \left(\sqrt{\frac{2r_2}{r_1+r_2}} - 1\right),\\
  \Delta v_2 &= \sqrt{\frac{GM_\odot}{r_2}}
               \left(1 - \sqrt{\frac{2r_1}{r_1+r_2}}\right),
\end{align}
转移时间（椭圆轨道的半个周期）：
\begin{equation}
  t_{\mathrm{霍曼}} = \pi\sqrt{\frac{(r_1+r_2)^3}{8\,GM_\odot}}.
\end{equation}
对于地球--火星霍曼转移（$r_1 = 1.000$\,AU，$r_2 = 1.524$\,AU）：
$\Delta v_1 \approx \SI{2.94}{\kilo\meter\per\second}$，
$\Delta v_2 \approx \SI{2.65}{\kilo\meter\per\second}$，
$t_{\mathrm{霍曼}} \approx \SI{259}{\day}$。

\begin{figure}[ht]
  \centering
  \begin{tikzpicture}[scale=0.95]
    \fill[yellow!90!orange] (0,0) circle (0.28cm);
    \node[below right, font=\footnotesize] at (0.15,-0.15) {太阳};
    \draw[blue!65, thick] (0,0) circle (2cm);
    \fill[blue!65] (2,0) circle (0.12cm);
    \node[above right, font=\footnotesize, blue!65] at (2.05,0.15) {地球 ($r_1$)};
    \draw[red!65, thick] (0,0) circle (3.05cm);
    \fill[red!65] (-3.05,0) circle (0.12cm);
    \node[above left, font=\footnotesize, red!65] at (-3.1,0.15) {火星 ($r_2$)};
    \draw[green!55!black, very thick, dashed,
          domain=0:180, samples=120]
      plot ({-0.525 + 2.525*cos(\x)}, {2.470*sin(\x)});
    \draw[-{Stealth[length=5pt]}, magenta!80!black, very thick]
      (2,0) -- (2,0.75);
    \node[right, font=\footnotesize, magenta!80!black] at (2.08,0.37)
      {$\Delta v_1$};
    \draw[-{Stealth[length=5pt]}, orange!80!black, very thick]
      (-3.05,0.75) -- (-3.05,0);
    \node[right, font=\footnotesize, orange!80!black] at (-2.98,0.37)
      {$\Delta v_2$};
    \draw[-{Stealth}, blue!65, semithick] (0.05,2) -- (-0.25,2);
    \draw[-{Stealth}, red!65, semithick] (0.05,3.05) -- (0.35,3.05);
    \node[green!50!black, font=\footnotesize, align=center] at (0,3.45)
      {转移轨道};
  \end{tikzpicture}
  \caption{地球（蓝色，内圆）到火星（红色，外圆）的霍曼转移。
    飞船在地球轨道施加冲量 $\Delta v_1$，历时约 259 天后在火星制动 $\Delta v_2$。}
  \label{fig:hohmann_zh}
\end{figure}

\subsection{引力辅助机动}

飞掠大质量行星可以从行星提取轨道动能并传递给飞船——代价几乎为零。在行星参考系中飞船速率不变，但方向偏转角为：
\begin{equation}
  \delta = 2\arcsin\!\left(\frac{1}{1 + r_p v_\infty^2/(GM_p)}\right).
\end{equation}
在日心参考系中的速度增量：
\begin{equation}
  \Delta v_\infty = 2v_p\sin\!\left(\frac{\delta}{2}\right).
\end{equation}
旅行者2号利用四次连续引力辅助，实现了太阳系逃逸。

\subsection{低能转移与星际高速公路}

在霍曼转移之外，\citet{Belbruno1987} 证明：日-地和日-月三体系统中
$L_1$、$L_2$ 平动点的不变流形在整个太阳系内编织出一张低能通道网络。
沿这些流形转移可以大幅降低 $\Delta v$，代价是更长的飞行时间，其本质
是利用三体问题对初始条件的敏感依赖性。支配方程为圆型限制性三体问题
方程：
\begin{align}
  \ddot{x} - 2\dot{y} &= x - \frac{(1-\mu)(x+\mu)}{r_1^3}
    - \frac{\mu(x-1+\mu)}{r_2^3}, \label{eq:cr3bp_x_zh}\\
  \ddot{y} + 2\dot{x} &= y - \frac{(1-\mu)\,y}{r_1^3}
    - \frac{\mu\,y}{r_2^3}, \label{eq:cr3bp_y_zh}
\end{align}
其中 $\mu = M_2/(M_1+M_2)$ 为质量参数，$r_{1,2}$ 为到各主天体的距离。
雅可比积分
\begin{equation}
  C_J = x^2 + y^2 + \frac{2(1-\mu)}{r_1} + \frac{2\mu}{r_2}
        - \dot{x}^2 - \dot{y}^2
  \label{eq:jacobi_zh}
\end{equation}
是圆型限制性三体问题唯一的守恒量，它定义了约束可能轨道的\emph{零速度面}。

% =============================================================================
\section{先进推进技术：超越化学火箭}
% =============================================================================

没有任何化学火箭能到达恒星。齐奥尔科夫斯基方程表明，所需速度增量对应的质量比根本无法实现。

\subsection{电推进与离子推进}

离子推进器将推进剂（通常为氙气）电离并通过静电加速，实现远高于化学发动机的喷气速度（$I_{\mathrm{sp}} \sim 3000$--$10000$\,s）。推力-功率关系为：
\begin{equation}
  F = \frac{2\eta_T\,P}{v_e},
\end{equation}
其中 $\eta_T \approx 0.65$ 为推进器效率。

\subsection{太阳帆}

距太阳 $d$ 处的辐射压强：
\begin{equation}
  P_{\mathrm{辐射}} = \frac{L_\odot}{4\pi d^2 c}.
\end{equation}
面质比为 $\sigma$ 的完全反射太阳帆所受加速度：
\begin{equation}
  a_{\mathrm{太阳帆}} = \frac{L_\odot\,\sigma}{2\pi d^2\,c}.
\end{equation}
当代超薄聚酯薄膜帆可达到
$\sigma \sim \SI{100}{\meter\squared\per\kilo\gram}$。

\subsection{核热推进}

核热火箭在反应堆中将推进剂（氢）加热到约 $\sim\SI{2700}{\kelvin}$，
可获得 $I_{\mathrm{sp}} \approx \SI{800}{}$--$\SI{1000}{\second}$——约为
最佳化学发动机的两倍。20世纪60年代的 NERVA 地面试验验证了这一技术。
$I_{\mathrm{sp}}$ 的提高使火星往返所需的推进剂质量减半。

\subsection{核脉冲推进（奥利安计划）}

奥利安计划提出在推进板后方引爆核弹产生冲量~\citep{Dyson1968}。理论比冲在 $10^4$ 至 $10^6$\,s 之间，可实现数周内的载人土星任务。

\subsection{反物质推进}

质子-反质子湮灭释放出质量所能提供的最大能量：
\begin{equation}
  E = mc^2 \quad \Rightarrow \quad
  I_{\mathrm{sp,最大}} = \frac{c}{g_0} \approx \SI{3.06e7}{\second}.
\end{equation}
反物质比例为 $\phi$ 时，有效喷气速度约为：
\begin{equation}
  v_e \approx c\sqrt{4\phi(1-\phi)},
\end{equation}
在 $\phi = 0.5$ 时达到峰值 $v_e = c$。

\subsection{激光推进（突破摄星计划）}

地面激光阵列功率 $P_L$，对质量 $m$、反射率 $R$ 的帆施加的加速度：
\begin{equation}
  a = \frac{(1+R)\,P_L}{mc}.
\end{equation}
突破摄星计划目标是将 \SI{1}{\gram} 的探测器加速到 $0.2c$，约20年抵达半人马座$\alpha$。

\subsection{星际冲压发动机}

\citet{Bussard1960} 指出，高速飞船不必携带全部推进剂，而可以从星际介质
中就地采集。磁性\emph{冲压收集器}会电离并汇聚周围的氢进入聚变室，使
飞船持续补充燃料。在速度 $v$ 下拦截的推进剂质量流量为
\begin{equation}
  \dot m = \rho\,A\,v,
  \label{eq:ram_flux_zh}
\end{equation}
其中 $\rho$ 为环境密度，$A$ 为有效收集面积。可用推力受聚变释放功率
$\dot m\,\varepsilon c^2$ 与所收集气体的阻力共同限制，因此极限速度为
\begin{equation}
  v_{\max} \approx c\sqrt{2\,\eta_f\,\varepsilon},
  \label{eq:ram_vmax_zh}
\end{equation}
$\varepsilon$ 为静质量转化为能量的比例，$\eta_f$ 为转化为定向推力的效率。
对质子-质子链（$\varepsilon \approx 0.007$）和理想耦合（$\eta_f = 1$），
$v_{\max} \approx 0.12c$。该构想优雅，但工程难度极高：收集器直径需达
数千公里，且在低速时所收集等离子体的轫致辐射损失超过聚变增益。因此，
冲压发动机更应视为长期愿景而非近期设计。

\begin{table}[ht]
  \centering
  \caption{推进技术比较。
           T.R.L.~=~技术成熟度（$1$\,=\,基本原理，$9$\,=\,飞行验证）。}
  \label{tab:propulsion_zh}
  \begin{tabularx}{\textwidth}{l r r c X}
    \toprule
    技术 & $I_{\mathrm{sp}}$ (s) & 推力 (N) & T.R.L. & 最佳用途 \\
    \midrule
    化学（LOX/LH$_2$）   & 450            & $\sim\!10^6$ & 9 & 低轨、行星际 \\
    煤油/液氧            & 310            & $\sim\!10^7$ & 9 & 发射 \\
    离子（NSTAR）        & 3\,100         & $0.1$        & 9 & 深空巡航 \\
    霍尔效应（SPT-140）  & 1\,600         & 5.6          & 8 & 地球同步轨道维持 \\
    太阳帆               & N/A            & $10^{-3}$    & 5 & 近日、小型探测器 \\
    核热                 & 900            & $\sim\!10^5$ & 4 & 火星、快速转移 \\
    核脉冲（奥利安）     & $10^4$--$10^6$ & $\sim\!10^7$ & 2 & 外太阳系、恒星际 \\
    聚变（D-T）          & $10^6$         & $10^3$       & 2 & 恒星际先驱 \\
    激光帆               & N/A            & 可变         & 2 & 飞往近邻恒星的纳米探测器 \\
    反物质               & $\sim\!3\!\times\!10^7$ & 可变 & 1 & 完整恒星际飞行 \\
    \bottomrule
  \end{tabularx}
\end{table}

% =============================================================================
\section{相对论性太空飞行}
% =============================================================================

\subsection{狭义相对论简介}

在接近光速时，牛顿力学失效。爱因斯坦狭义相对论（1905）~\citep{Einstein1905}引入洛伦兹因子：
\begin{equation}
  \gamma \equiv \frac{1}{\sqrt{1-\beta^2}}, \qquad \beta \equiv v/c.
\end{equation}
其主要推论包括：时间膨胀 ($\Delta t' = \gamma\,\Delta\tau$)，长度收缩 ($L = L_0/\gamma$)，以及能量-动量关系：
\begin{equation}
  E^2 = (pc)^2 + (mc^2)^2.
\end{equation}

\begin{figure}[ht]
  \centering
  \begin{tikzpicture}
    \begin{axis}[
      width=11cm, height=7cm,
      xlabel={$\beta = v/c$},
      ylabel={洛伦兹因子 $\gamma$},
      xmin=0, xmax=1, ymin=1, ymax=12,
      grid=both,
      grid style={line width=0.3pt, draw=gray!25},
      major grid style={line width=0.5pt, draw=gray!50},
      tick label style={font=\small},
      label style={font=\small},
      legend pos=north west, legend style={font=\small},
    ]
      \addplot[blue!80!black, very thick, domain=0:0.997, samples=300]
        {1/sqrt(1-x*x)};
      \addlegendentry{$\gamma=(1-\beta^2)^{-1/2}$}
      \addplot[only marks, mark=*, mark size=2pt, red!70]
        coordinates {(0.5,1.155)(0.8,1.667)(0.9,2.294)(0.98,5.025)};
    \end{axis}
  \end{tikzpicture}
  \caption{洛伦兹因子 $\gamma$ 作为 $\beta = v/c$ 的函数。
    当 $\beta \to 1$ 时，$\gamma$ 急剧上升：$\beta = 0.99$ 时，飞船上的时钟比静止时慢7倍以上。}
  \label{fig:lorentz_zh}
\end{figure}

\subsection{相对论火箭方程}

齐奥尔科夫斯基方程的相对论推广~\citep{Ackeret1946}为：
\begin{equation}
  \boxed{\frac{m_0}{m_f} = \left(\frac{1+\beta_f}{1-\beta_f}\right)^{c/(2v_e)},}
\end{equation}
或用快度 $\phi = \atanh(\beta)$ 表示：
\begin{equation}
  \frac{m_0}{m_f} = \exp\!\left(\frac{c}{v_e}\,\atanh(\beta_f)\right).
\end{equation}
对于完美反物质推进器（$v_e = c$），末速 $\beta_f = 0.99$ 时：$m_0/m_f \approx 14$——这是一个可实现的比值。

\subsection{匀加速轨迹}

在恒定本征加速度 $a$ 下，飞船的轨迹（林德勒双曲线）满足：
\begin{align}
  x(\tau) &= \frac{c^2}{a}\!\left(\cosh\frac{a\tau}{c} - 1\right),\\
  t(\tau)  &= \frac{c}{a}\sinh\frac{a\tau}{c},\\
  v(\tau)  &= c\tanh\frac{a\tau}{c},
\end{align}
其中 $\tau$ 为旅行者的固有时间。对于坐标距离为 $D$ 的单程航行（前半程
加速、后半程减速），所经历的本征时间为
\begin{equation}
  \tau = 2\,\frac{c}{a}\,\cosh^{-1}\!\!\left(\frac{aD}{2c^2}+1\right),
  \label{eq:oneway_proper_zh}
\end{equation}
而地球上测得的坐标时间为
\begin{equation}
  t = 2\,\frac{c}{a}\sinh\!\left(\frac{a\tau}{2c}\right)
    = 2\,\frac{c}{a}\sqrt{\left(\frac{aD}{2c^{2}}+1\right)^{2}-1}.
  \label{eq:oneway_coord_zh}
\end{equation}
往返航行则将两者加倍。当 $a = g = \SI{9.81}{\meter\per\second\squared}$
（舒适的 1~$g$ 加速度）时：

\begin{table}[ht]
  \centering
  \caption{1~$g$ 连续加速下飞往若干目的地的单程时间（$a=g$）。
   坐标时间 $t$ 由地球上的静止观测者测量，$\tau$ 为旅行者的固有时间。
   往返时间为此值的两倍。}
  \label{tab:rel_trips_zh}
  \begin{tabularx}{\textwidth}{X r r r}
    \toprule
    目的地 & $D$ & $t$（坐标） & $\tau$（固有） \\
    \midrule
    火星（平均）        & \SI{0.524}{\AU} & 2.1 天        & 2.1 天        \\
    木星                & \SI{5.2}{\AU}   & 6.5 天        & 6.5 天        \\
    冥王星              & \SI{39.5}{\AU}  & 18.0 天       & 18.0 天       \\
    比邻星              & \SI{4.24}{\ly}  & 5.9 年        & 3.5 年        \\
    银河系中心          & \SI{26000}{\ly} & 26\,000 年    & 19.8 年       \\
    仙女座星系          & \SI{2.5e6}{\ly} & 250 万年      & 28.6 年       \\
    \bottomrule
  \end{tabularx}
\end{table}

时间膨胀效应十分深远。以 1~$g$ 飞往比邻星，飞船时间仅需 3.5 年，而地球
上已过去 5.9 年；往返则使船员老化 7.1 年，地球过去 11.7 年。飞往银河系
中心只会让船员老化 39.8 年——但地球上将流逝超过 52000 年。宇宙并非仅仅
可以航行：在 1~$g$ 下，它对旅行者而言小得令人惊讶。

\subsection{能量需求}

相对论性动能为：
\begin{equation}
  E_{\mathrm{动}} = (\gamma_f - 1)\,m_f\,c^2.
\end{equation}
将 $10^6$\,kg 有效载荷加速到 $\beta=0.2$，此时
$\gamma_f = (1-0.04)^{-1/2} \approx 1.0206$，则
\begin{equation*}
  E = 0.0206 \times 10^6 \times (2.998\times10^{8})^2
    \approx \SI{1.85e21}{\joule} \approx \SI{5.1e5}{\tera\watt\hour}.
\end{equation*}
全球一次能源年消耗量约为 \SI{6.2e20}{\joule}，因此将单个 1000 吨有效载荷
加速到 $0.2c$ 约需消耗当今世界三年的能源供应。这需要能够持续输出
$10^{17}$--$10^{19}$\,W 的先进聚变或反物质反应堆——完全处于卡尔达肖夫
II 型文明的能力范围之内。

\subsection{星际尘埃、辐射与防护}

在 $\beta \sim 0.2$--$0.9$ 时，星际介质——尽管比任何地面实验室所能制造的
真空还要稀薄，平均密度约为每立方厘米一个原子——也会成为工程风险。静质
量为 $m_g$ 的尘埃颗粒被正面撞击时释放的动能为
\begin{equation}
  E_g = (\gamma - 1)\,m_g\,c^2,
  \label{eq:dust_impact_zh}
\end{equation}
因此 \SI{e-12}{\kilo\gram} 的颗粒（微米尺度的粒子）在 $\beta = 0.2$ 时释放
约 \SI{1.9e3}{\joule}，在 $\beta = 0.9$ 时释放 \SI{1.2e5}{\joule}——相当于
重型炮弹的炮口动能。尘埃颗粒的空间密度很低，但跨越数光年的航行所扫过的
体积如此之大，以至于薄的前向防护盾或磁/等离子体偏转器是必需的。星际氢
本身也会造成持续侵蚀，在最高速度下还会由散裂产物产生贯穿辐射。任何可信
的相对论性星舰设计，都必须将防护与避碰视为首要约束，而非事后补充
\citep{Dyson1968}。

% =============================================================================
\section{文明数量的统计论证}
% =============================================================================

\subsection{德雷克方程}

银河系中目前有多少技术活跃的文明能够进行星际通信？弗兰克·德雷克（1961）~\citep{Drake1961}提出了基础估算：
\begin{equation}
  \boxed{N = R_* \cdot f_p \cdot n_e \cdot f_l \cdot f_i \cdot f_c \cdot L,}
\end{equation}
其中：$R_*$ 为银河系恒星形成率（年$^{-1}$），$f_p$ 为有行星系统的恒星比例，$n_e$ 为每个行星系统中宜居带行星的平均数量，$f_l$ 为出现生命的比例，$f_i$ 为发展出智慧生命的比例，$f_c$ 为能够进行可探测星际通信的比例，$L$ 为此类文明的平均寿命（年）。

表~\ref{tab:drake_zh} 汇总了各参数的当前估计。

\begin{table}[ht]
  \centering
  \caption{德雷克方程参数：保守、乐观与最佳当前估计。
    所得 $N$ 跨越十一个数量级以上。}
  \label{tab:drake_zh}
  \begin{tabularx}{\textwidth}{X r r r}
    \toprule
    参数 & 保守 & 乐观 & 最佳估计 \\
    \midrule
    $R_*$（年$^{-1}$） & 1    & 3          & 3    \\
    $f_p$              & 0.5  & 1.0        & 0.9  \\
    $n_e$              & 0.01 & 1.0        & 0.4  \\
    $f_l$              & $10^{-6}$ & 1.0  & 0.1  \\
    $f_i$              & $10^{-6}$ & 1.0  & 0.01 \\
    $f_c$              & 0.01 & 0.5        & 0.1  \\
    $L$（年）          & 100  & $10^9$     & $10^4$ \\
    \midrule
    $N$       & $\sim 10^{-10}$ & $\sim 10^8$ & $\sim 1200$ \\
    \bottomrule
  \end{tabularx}
\end{table}

当前估计结果从 $N \approx 0$（悲观）到 $N \sim 10^8$（乐观），合理的中心值约为 $N \sim 1200$。

\subsection{贝叶斯重构}

将通信文明的空间密度建模为泊松过程，密度为 $\lambda_c$，在体积 $V$ 内至少存在一个文明的概率为：
\begin{equation}
  P(N \geq 1) = 1 - \ee^{-\lambda_c V}.
\end{equation}

\subsection{费米悖论及其解释}

若 $N$ 很大，费米的问题便无法回避~\citep{Hart1975}：\textit{"他们在哪里？"}主要解释可归为三类：
\begin{enumerate}[noitemsep]
  \item \textbf{稀有性。} 生命或智慧极为罕见（"稀有地球"假说~\citep{WardBrownlee2000}），$N \ll 1$。
  \item \textbf{大过滤器。} 存在某个极低概率的进化跃迁，或在我们之前，或在我们之后~\citep{Hanson1998}。
  \item \textbf{通信差异。} 先进文明存在，但以我们未曾关注的方式通信（动物园假说、黑暗森林理论~\citep{Liu2008}）。
\end{enumerate}

\subsection{卡尔达肖夫等级}
\label{sec:kardashev_zh}

尼古拉·卡尔达肖夫（1964）按能源消耗将文明分类：
\begin{itemize}[noitemsep]
  \item \textbf{I型：} 利用母星球全部能量，$P_{\mathrm{I}} \sim \SI{e16}{\watt}$。
  \item \textbf{II型：} 驾驭恒星全部能量，$P_{\mathrm{II}} \sim \SI{4e26}{\watt}$。
  \item \textbf{III型：} 掌控整个星系能量，$P_{\mathrm{III}} \sim \SI{e37}{\watt}$。
\end{itemize}
连续化公式~\citep{Sagan1973}：$K = (\log_{10} P_W - 6)/10$。
2026年人类消耗约 $\SI{2e13}{\watt}$，对应 $K \approx 0.73$，接近但尚未达到I型。

% =============================================================================
\section{跨越星际距离的通信}
% =============================================================================

\subsection{电磁通信}

发射机功率为 $P_t$、增益为 $G_t$，信号在距离 $d$ 处的弗里斯传输方程：
\begin{equation}
  P_r = P_t\,G_t\,G_r\!\left(\frac{\lambda}{4\pi d}\right)^2.
\end{equation}
带宽 $B$、系统噪声温度 $T_{\rm sys}$ 下的信噪比：
\begin{equation}
  \mathrm{SNR} = \frac{P_t\,G_t\,A_r}{4\pi d^2\,\kB\,T_{\rm sys}\,B}.
\end{equation}

\subsection{宇宙"聚水坑"}

氢线（1420.4\,MHz）与羟基线（1612--1720\,MHz）之间的频谱窗口在宇宙背景中相对安静，被称为宇宙"聚水坑"，是信标信号的天然集合点。

\subsection{信道容量估算}

香农信道容量定理~\citep{Shannon1948}：
\begin{equation}
  C = B\log_2\!\left(1 + \mathrm{SNR}\right).
\end{equation}
一个功率为 \SI{10}{\tera\watt} 的发射机，在10光年外用100米口径射电望远镜完全可检测——只要知道搜索方向和频率。

% =============================================================================
\section{首次接触的通用语言}
% =============================================================================

\subsection{数学的普适性}

任何足够先进的文明必定已经发现了相同的数学真理：2、3、5、7、11的素性；$\pi$ 的值；欧几里得几何的结构；以及自然界的物理常数。数学不是人类的发明，而是对物理现实中固有模式的发现——是任何掌握技术的智慧生命都能解读的唯一框架~\citep{Penrose1989}。

\subsection{林科斯方案与素数编码}

汉斯·弗罗伊登塔尔（1960）~\citep{Freudenthal1960}提出从零开始构建一种星际语言（林科斯）。编码层次如下：
\begin{enumerate}[noitemsep]
  \item \textbf{数字。} 以素数数量的脉冲串（2、3、5、7、11、……）建立计数。
  \item \textbf{算术。} 以脉冲串形式传递 $2 + 3 = 5$ 等算术关系。
  \item \textbf{逻辑。} 布尔常数与命题算符。
  \item \textbf{关系。} 等式、序关系、集合成员关系。
  \item \textbf{物理参照量。} 氢原子质量、无量纲常数。
\end{enumerate}
最重要的无量纲自然常数——精细结构常数 $\alpha$ 和质子-电子质量比 $\mu$——是普适信标：
\begin{equation}
  \alpha = \frac{e^2}{4\pi\varepsilon_0\,\hbar c} \approx \frac{1}{137.036},
  \qquad
  \mu = \frac{m_p}{m_e} \approx 1836.15.
\end{equation}

\subsection{二维像素编码}

将总比特数 $N = p \times q$（$p$、$q$ 为不同的素数）映射为 $p \times q$ 像素网格，接收方可通过唯一素因数分解重构图像。1974年阿雷西博信息正是采用了这一技术：$1679 = 73 \times 23$ 比特，编码了一幅 $73 \times 23$ 像素的图像，包含太阳系、DNA结构和人类形象。

% =============================================================================
\section{建立星际关系：穿越虚空的外交}
% =============================================================================

\subsection{首次接触的博弈论}

接触问题可建模为协调博弈。设文明A和B各自选择策略 $s \in \{S, \varnothing\}$（发送信号或保持沉默）。收益矩阵为：
\begin{equation}
\begin{pmatrix}
  & S & \varnothing \\
  S & (b-c,\; b-c) & (-c,\; b) \\
  \varnothing & (b,\; -c) & (0,\; 0)
\end{pmatrix}
\end{equation}
若 $b > c$，则相互发信 $(S, S)$ 是纳什均衡。若 $b < 0$（接触有危险），则沉默成为双方的占优策略——即黑暗森林的博弈论解释~\citep{Liu2008}。

\subsection{迭代接触与信任建立}

在迭代囚徒困境中，"以牙还牙"策略（第一轮合作，此后镜像对方上一轮行动）被证明具有极强鲁棒性。在 $T$ 轮博弈中的期望收益：
\begin{equation}
  U(T) = b + (T-1)\left[b\,p_{\rm 合作} - c\,p_{\rm 背叛}\right].
\end{equation}
对于长寿的文明（$T \gg 1$），只要 $b > c$，合作占主导——这是一个令人鼓舞的结论。

\subsection{接触协议}

任何文明化的接触体系都应包含：(i) 以数学为媒介的意图声明；(ii) 考虑光年级往返时延的时间耐心；(iii) 对未知文明的全向广播；(iv) 由机构（而非个人）维系这场跨越数代人的跨星际通信。

\subsection{接触的伦理}

若我们探测到外星信号，三种伦理框架给出不同答案：
（一）功利主义——若预期信息收益巨大且风险可控，则积极回应是合理的；
（二）预防原则——向未知文明暴露我们的位置存在潜在生存风险，宜保持沉默；
（三）世界主义——所有智慧生命共享一个道德共同体，拒绝回应是一种狭隘的孤立主义。
目前，国际天文学联合会SETI委员会建议在国际社会就回应协议达成共识之前，仅被动接收信号。

% =============================================================================
\section{迈向银河文明}
% =============================================================================

\subsection{信息、能量与复杂性}

朗道尔原理表明，物理信息是有代价的：
\begin{equation}
  \Delta Q_{\min} \geq \kB T \ln 2 \approx \SI{2.87e-21}{\joule}
  \quad (T = \SI{300}{\kelvin}).
\end{equation}
文明的信息处理能力与其能源预算成比例。从卡尔达肖夫I型到II型到III型的跃迁，对应计算能力约 $10^{10}$ 倍的提升。

\subsection{文明的演化轨迹}

若文明能够可靠地达到II型（约有 $10^{10}$ 年的恒星演化时间可供利用），
那么银河系早已有足够的时间被反复殖民。而事实并非如此——或看起来并非
如此——这一观察要么是支持``大过滤器在我们前方''的最强证据，要么表明
扩张中的文明会收敛到稳定、低足迹的形态（\emph{超越假说}：先进智能
向内发展，追求不断增长的复杂性与微型化，而非向外进行领土扩张）。

\subsection{文明的长远轨迹}

"恒星纪元"将持续约 $\sim 10^{14}$ 年；红矮星和棕矮星将在太阳成为白矮星之后很久，仍能支撑宜居行星。一个在未来数千年中存活下来的文明将发现，宇宙以当前的恒星形成率衡量，尚未度过其产星寿命的一半。遥远未来的道德重量——文明延续所能孕育的潜在存在——是将生命扩展至宇宙视为最高文明优先事项的压倒性理由。

% =============================================================================
\section{结论}
% =============================================================================

我们已经从行星运动的活力公式到应用于星际信标的香农信道容量定理，梳理了一条连贯的数学主线——这条线索连接了轨道力学的物理学、火箭推进的热力学、时空的几何学与语言的信息论。

分析得出了确定性的结论：

\begin{enumerate}[noitemsep]
  \item 化学推进从根本上无法实现恒星际飞行；所需的最低推进技术为聚变推进、反物质推进或激光推进（用于小型探测器）。
  \item 相对论性飞行到最近恒星在物理上可行，且对飞船上的人而言可在有生之年内完成，但所需的能源预算要求达到卡尔达肖夫I型甚至II型文明。
  \item 德雷克方程既与 $N \gg 1$ 相容，也与 $N \approx 0$ 相容；未来的生物特征巡天（詹姆斯·韦伯空间望远镜等）将大幅缩小这一不确定范围。
  \item 千兆瓦至太瓦级功率的电磁通信，在数十光年外用近未来射电和光学望远镜完全可探测；主要障碍在于指向、频率和时机。
  \item 数学是首次接触唯一中性的语言：素数、普适常数和几何编码是其共同基础。
  \item 当文明寿命漫长且接触是迭代的，博弈论有利于星际合作；在大多数合理参数选择下，相互通信的长期期望值为正。
\end{enumerate}

星辰并非遥不可及。时间、能量和勇气将我们与它们分隔——而这三者，人类在其短暂的宇宙历史中，已经证明自己丰厚地拥有。问题不是我们\textit{是否}会到达那里，而是\textit{何时}，以及我们到达时会\textit{遇见谁}。

% =============================================================================
% 附录
% =============================================================================
\appendix

\section{相对论火箭方程的推导}

在火箭参考系中以速度 $v_e$ 喷射推进剂时，四动量守恒给出：
\begin{equation}
  \int_0^{\beta_f} \frac{c\,\dd\beta'}{(1-\beta'^2)(1 - v_e\beta'/c^2)}
  = -v_e \ln\frac{m_f}{m_0}.
\end{equation}
对光子推进（$v_e = c$），积分值为 $c\,\atanh(\beta_f)$，直接给出相对论火箭方程。

\section{基本物理常数数值}

\begin{tabular}{ll}
  \toprule
  常数 & 数值 \\
  \midrule
  光速 & $c = \SI{2.998e8}{\meter\per\second}$ \\
  引力常数 & $G = \SI{6.674e-11}{\meter\cubed\per\kilo\gram\per\second\squared}$ \\
  玻尔兹曼常数 & $\kB = \SI{1.381e-23}{\joule\per\kelvin}$ \\
  太阳质量 & $\Msun = \SI{1.989e30}{\kilo\gram}$ \\
  太阳光度 & $L_\odot = \SI{3.828e26}{\watt}$ \\
  天文单位 & $\SI{1}{\AU} = \SI{1.496e11}{\meter}$ \\
  光年 & $\SI{1}{\ly} = \SI{9.461e15}{\meter}$ \\
  Hubble 常数 & $H_0 = 67.4$\,km\,s$^{-1}$\,Mpc$^{-1}$ \\
  \bottomrule
\end{tabular}

% =============================================================================
% 参考文献
% =============================================================================
\bibliographystyle{plainnat}
\bibliography{references}

\end{document}
